\section{Experiments}
\label{sec:experiments}

\subsection{Setup}
\paragraph{Benchmark.} We evaluate on PIE-Bench \citep{ju2024pnpinversion}, 700 real images with source/target
prompts and edit-region masks over 10 edit types. Because \method{} targets non-rigid edits, we additionally
report the \emph{change pose} subset (40 images).

\paragraph{Metrics.} Following PIE-Bench, structure distance \citep{tumanyan2022splicing} measured with DINO
features \citep{caron2021dino}; PSNR, LPIPS \citep{zhang2018lpips}, MSE and SSIM on the region outside the
edit mask (background preservation); and CLIP similarity \citep{radford2021clip} between the edited image
(whole / edited region) and the target prompt (edit fidelity).

\paragraph{Baselines and implementation.} All methods use Stable Diffusion v1.4 \citep{rombach2022ldm},
DDIM with 50 steps and classifier-free guidance 7.5, $S=4$, $L=10$, and the same seed.
We compare against (i) \emph{fixed seed}: the inverted noise denoised with the target prompt and no attention
control; (ii) DDIM inversion + \method{} using the official implementation; and (iii) Direct Inversion
\citep{ju2024pnpinversion} + \method{}. \todo{GPU type, runtime.}

\subsection{Main Results}
% Placeholder - overwritten by `python src/summarize.py` after the baseline runs.
\begin{table*}[t]
\centering
\small
\begin{tabular}{lccccccc}
\toprule
Method & Struct.$\downarrow$ & PSNR$\uparrow$ & LPIPS$\downarrow$ & MSE$\downarrow$ & SSIM$\uparrow$ & CLIP-W$\uparrow$ & CLIP-E$\uparrow$ \\
\midrule
SD fixed seed (no control) & -- & -- & -- & -- & -- & -- & -- \\
DDIM + MasaCtrl & -- & -- & -- & -- & -- & -- & -- \\
DirectInv. + MasaCtrl & -- & -- & -- & -- & -- & -- & -- \\
\midrule
\textit{DDIM+MasaCtrl (PIE-Bench paper)} & 28.38 & 22.17 & 106.62 & 86.97 & 79.67 & 23.96 & 21.16 \\
\textit{DirectInv.+MasaCtrl (PIE-Bench paper)} & 24.70 & 22.64 & 87.94 & 81.09 & 81.33 & 24.38 & 21.35 \\
\bottomrule
\end{tabular}
\caption{Baseline results on PIE-Bench (700 images). Struct.\ $\times10^3$, LPIPS $\times10^3$, MSE $\times10^4$, SSIM $\times10^2$; PSNR/LPIPS/MSE/SSIM computed on the unedited region, CLIP-W/E on the whole/edited region. \textcolor{red}{Placeholder until our runs finish.}}
\label{tab:baseline}
\end{table*}

\begin{table*}[t]
\centering
\small
\begin{tabular}{lccccccc}
\toprule
Method & Struct.$\downarrow$ & PSNR$\uparrow$ & LPIPS$\downarrow$ & MSE$\downarrow$ & SSIM$\uparrow$ & CLIP-W$\uparrow$ & CLIP-E$\uparrow$ \\
\midrule
SD fixed seed (no control) & -- & -- & -- & -- & -- & -- & -- \\
DDIM + MasaCtrl & -- & -- & -- & -- & -- & -- & -- \\
DirectInv. + MasaCtrl & -- & -- & -- & -- & -- & -- & -- \\
\bottomrule
\end{tabular}
\caption{Baseline results on the PIE-Bench \emph{change pose} subset (40 images), the non-rigid setting MasaCtrl targets.}
\label{tab:baseline_pose}
\end{table*}

\todo{Discuss Table~\ref{tab:baseline} and Table~\ref{tab:baseline_pose}; sanity-check our DDIM+\method{}
reproduction against the published PIE-Bench numbers.}

\subsection{Qualitative Results}
\todo{Figure: source $|$ fixed seed $|$ \method{} $|$ \ours{} on 4--5 pose edits.}

\subsection{Ablation Study}
\todo{Effect of $S$ and $L$; components of \ours{}.}
